\newpage
\section{Discussion}

APEIRON-Core demonstrates that a unified, low-overhead hardware abstraction layer spanning CPU, GPU, FPGA, and quantum backends is feasible and practical. The experimental results confirm that the abstraction tax is negligible for realistic workloads, that significant speedups are achievable on specialized hardware, and that the error-handling and energy-aware mechanisms operate correctly. Critically, the system meets the demands of non-anthropocentric AI: it allows higher cognitive layers to express computations in a hardware-agnostic manner while dynamically optimizing for performance, reliability, and thermodynamic cost.

Regarding the ``breadth vs. depth'' concern: We acknowledge that each backend (especially FPGA and quantum) could merit a full paper on its own. Our contribution is not to advance the state of the art in any single hardware domain, but rather to provide a unifying framework that makes these disparate resources accessible to a common AI substrate. The depth is present in the careful handling of cross-backend impedance mismatches (e.g., fixed-point conversion, SWAP test correction) and in the integrated observability and fallback mechanisms. We view this work as enabling future research that can leverage any computational matter without rewriting the core knowledge algorithms.

\subsection{Limitations}

\begin{itemize}
    \item FPGA Portability: The current FPGA overlay is hand-optimized for the ZCU104. A HLS-based generator that can target different Zynq devices is under development.
    \item Quantum Scalability: Statevector simulation limits the quantum backend to $\sim$10 qubits on local hardware. We plan to integrate circuit-cutting techniques to simulate larger systems classically and to more fully utilize IBM Quantum's dynamic circuits for hardware execution.
    \item Energy Model Fidelity: The carbon-aware throttling currently uses coarse-grained grid intensity; a per-workload energy model based on hardware performance counters would improve accuracy and enable finer-grained decisions.
    \item Benchmark Representativeness: The workloads used are synthetic but derived from APEIRON's internal loops. A full-scale APEIRON simulation (with $10^5$ fields evolving over $10^6$ timesteps) would provide a more definitive validation, but is currently computationally prohibitive.
\end{itemize}

\subsection{Future Work}

\begin{itemize}
    \item Extend the FPGA overlay to support higher-order interference detection and on-chip pattern clustering.
    \item Implement a neuromorphic backend for Intel Loihi 2, leveraging Lava for event-driven field dynamics.
    \item Develop a workload-aware energy profiler that uses RAPL and NVML to build per-operation energy models, enabling predictive carbon-aware scheduling.
    \item Release the benchmark suite and FPGA bitstream as open-source artifacts to facilitate reproducibility and community extension.
\end{itemize}